\documentclass[aspectratio=169, t, 8pt]{beamer}

%------configurando acentos-----------
\usepackage[utf8]{inputenc}
\usepackage[brazil]{babel}

\usepackage{graphicx}

%-------Estilos Beamer---------------
\usetheme{metropolis}
%\usecolortheme{crane}
%coment
%\usefonttheme{serif}
\linespread{1.2}

% Personalização de Cores (Exemplo: Minimal Teal)
% Cores Base
\definecolor{BgCanvas}{HTML}{F9F9FB}      % Off-white muito suave
\definecolor{MainText}{HTML}{2B2D42}      % Grafite escuro para alta legibilidade
\definecolor{PrimaryAccent}{HTML}{4A6FA5} % Az
\definecolor{SunCoastFg}{HTML}{F2A154}
% --- CONFIGURAÇÕES DE TEMA BEAMER ---
%\usetheme{metropolis}

\setbeamertemplate{itemize item}{\textcolor{MainText}{$\bullet$}}
\setbeamertemplate{itemize subitem}{\textcolor{MainText}{\scriptsize$\circ$}}
\setbeamertemplate{itemize 3rditem}{\textcolor{MainText}{\tiny$\blacksquare$}}

\setbeamercolor{background canvas}{bg=BgCanvas}
\setbeamercolor{normal text}{fg=MainText}
\setbeamercolor{frametitle}{bg=SunCoastFg, fg=BgCanvas}
\setbeamercolor{title}{fg=MainText}
\setbeamercolor{subtitle}{fg=PrimaryAccent}
\setbeamercolor{structure}{fg=PrimaryAccent}

% Remove a linha amarela padrão do Metropolis sob o título
\setbeamertemplate{title separator}{}
\setbeamertemplate{progress bar in section page}{}


%--------Dados do Documento
\title{Matemática: Lista de Exercícios}
\author{prof. Eduardo}
\institute{Pré-Universitário Oficina do Saber / UFF}
\date{\today}


%-------- Pacotes Matemática -------------
\usepackage{amsmath}
\usepackage{amsthm}
\usepackage{amsfonts}
\usepackage{amssymb}
\usepackage{xfrac}
\usepackage{thmtools}

%-------- Configuração Listas---------
% \usepackage{enumitem}
\usepackage[inline]{enumitem}

\setcounter{tocdepth}{2}

\newcounter{questaovest}
\newcommand{\vest}[1]{\stepcounter{questaovest}\item[\textbf{Questão \arabic{questaovest}} (#1):]}
\setlist[enumerate]{
    align=left,
    leftmargin=1.8em,
    itemindent=1.8em
    }

\setlist[enumerate,2]{
    label=\alph*.,
    font=\bfseries,
    itemsep=0.1pt,
    labelsep=0em,
    parsep=0pt,
    leftmargin=0em,
    }

%------------- Atalho para imagens ------------
\graphicspath{{../../src/img/semana08/}}


\begin{document}

\begin{frame} %capa
        \vspace*{\fill}
            
            \begin{columns}[c]
                \begin{column}{0.5\textwidth}
                    {\usebeamerfont{title}\usebeamercolor[fg]{title}\inserttitle}\vspace{0.2cm}
                    
                    {\usebeamerfont{subtitle}\usebeamercolor[fg]{subtitle}\insertsubtitle}\vspace{0.5cm}

                {\usebeamerfont{author}\insertauthor}\vspace{0.3cm}

                {\usebeamerfont{institute}\insertinstitute}\vspace{0.3cm}

                {\usebeamerfont{date}\insertdate}
            \end{column}
            
            
            \begin{column}{0.4\textwidth}
                \centering
                \includegraphics[width=0.4\textwidth]{../logo_ofSaber.png}

                \vspace{.4cm}
                \vspace*{\fill}
        \begin{columns}[c]
            \begin{column}{.5\textwidth}
                \centering\includegraphics[width=.5\textwidth]{../Proex02c.png}
            \end{column}
            \begin{column}{.5\textwidth}
                \centering\includegraphics[width=0.4\textwidth]{../Logo_UFF.png}
            \end{column}
        \end{columns}
        \vspace*{\fill}
            \end{column}
        \end{columns}
    \vspace*{\fill}
    \end{frame}
    
        % {
        %     \usebackgroundtemplate{
        %     \includegraphics[width=\paperwidth]{slide_02_contracapa.png}}
        %     \begin{frame}[plain]                
        %     \end{frame}
        % }


    \begin{frame} %Questão 01-C
        \frametitle{Razão e Proporção}
        \vspace{-1.2em}
        \begin{enumerate}
            \begin{columns}
                \begin{column}{.9\textwidth}
                    \vest{ENEM-2020} Um motociclista planeja realizar uma viagem cujo destino fica a 500 km de sua casa. Sua moto consome 5 litros de gasolina para cada 100 km rodados, e o tanque da moto tem capacidade para 22 litros. Pelo mapa, observou que no trajeto da viagem o último posto disponível para reabastecimento, chamado Estrela, fica a 80 km do seu destino.
                    
                    Ele pretende partir com o tanque da moto cheio e planeja fazer somente duas paradas para reabastecimento, uma na ida e outra na volta, ambas no posto Estrela. No reabastecimento para a viagem de ida, deve considerar também combustível suficiente para se deslocar por 200 km no seu destino.
                    
                    A quantidade mínima de combustível, em litro, que esse motociclista deve reabastecer no posto Estrela na viagem de ida, que seja suficiente para fazer o segundo reabastecimento, é   

                    \vspace{1em}
                            \raggedleft
                            \begin{minipage}{\textwidth}
                                \begin{enumerate*}[itemjoin={\hfill}]
                                    \item 13.
                                    \item 14.
                                    \item 17.
                                    \item 18.
                                    \item 21.
                                    %gabarito: C
                                \end{enumerate*}
                            \end{minipage}

                \end{column}
            \end{columns}
        \end{enumerate}
    \end{frame}

    \begin{frame} %Questão 02-C
        \frametitle{Geometria Espacial}
        \begin{columns}
            \begin{column}{.9\textwidth}
                \begin{enumerate}
                    \vest{UEPB-2013} Um reservatório em forma de cubo, cuja diagonal mede $2\sqrt{3}$ m, tem
                    \vspace{1.2em}
                    \begin{columns}[T, totalwidth=.8\textwidth]
                        \begin{column}{.35\textwidth}
                            \begin{enumerate}
                                \item 4000 litros
                                \item 6000 litros
                                \item 8000 litros
                                \item 2000 litros
                                \item 1000 litros
                            \end{enumerate}
                        \end{column}
                    \end{columns}
                \end{enumerate}
            \end{column}
        \end{columns}
    \end{frame}


    \begin{frame}{Contagem} %E
        \begin{columns}[totalwidth=.9\textwidth]
            \begin{column}{\textwidth}
                \begin{enumerate}
                    \vest{IFAL-2018} Em uma civilização antiga, o alfabeto tinha apenas três letras. Na linguagem dessa civilização, as palavras tinham de uma a quatro letras. Quantas palavras poderiam existir na linguagem dessa civilização?

                    \begin{enumerate}
                        \item 4
                        \item 12
                        \item 16
                        \item 40
                        \item 120
                    \end{enumerate}
                \end{enumerate}
            \end{column}
        \end{columns}
    \end{frame}

    \begin{frame}{Função Afim} %D
        \begin{columns}[totalwidth=.9\textwidth]
            \begin{column}{\textwidth}
                \begin{enumerate}
                    \vest{UEMA} Uma pequena empresa que fabrica camisetas verificou que o lucro obtido com a venda de seus produtos obedece à função $L(x) = 75x – 3000$, sendo $L(x)$ o lucro em reais e x o número de camisetas vendidas, para $40 < x \leq 120$.
                    
                    Para que o lucro da empresa chegue a R\$ 4.000,00, o menor número de camisetas a serem vendidas é:

                    \vspace{1.2em}

                    \begin{enumerate}
                        \item 97.
                        \item 96.
                        \item 95.
                        \item 94.
                        \item 93.
                    \end{enumerate}
                \end{enumerate}
            \end{column}
        \end{columns}
    \end{frame}

    \begin{frame}{Geometria Espacial}
        \begin{columns}[totalwidth=.9\textwidth]
            \begin{column}{\textwidth}
                \begin{enumerate}
                    \vest{Toda Matéria} Qual a área lateral de um cilindro reto que possui $502,4 \text{ cm}^3$ de volume e diâmetro $8 \text{ cm}$.
                    
                    Dado: $\pi = 3,14$.

                    \vspace{1.2em}

                    \begin{enumerate}
                        \item $355,10 \text{ cm}^2$
                        \item $251,20 \text{ cm}^2$
                        \item $125,51 \text{ cm}^2$
                        \item $375,30 \text{ cm}^2$
                        \item $91,45 \text{ cm}^2$
                    \end{enumerate}
                \end{enumerate}
            \end{column}
        \end{columns}
    \end{frame}

    % \begin{frame} %Questão 03-A
    %     \frametitle{Volume}
    %     \begin{columns}[T, totalwidth=.9\textwidth]
    %         \begin{column}{\textwidth}
    %             \begin{enumerate}
    %                 \vest{ENEM-2017} Para as Olimpíadas de 2012, a piscina principal do Centro Aquático de Londres, medindo 50 metros de comprimento, foi remodelada para ajudar os atletas a melhorar suas marcas. Observe duas das melhorias:
    %                 \begin{columns}[T, totalwidth=.9\textwidth]
    %                     \begin{column}{.6\textwidth}
    %                         \centering
    %                         \includegraphics[width=\textwidth]{monitoria_questao03.png}
    %                     \end{column}
    %                     \begin{column}{.35\textwidth}
    %                         A capacidade da piscina em destaque, em metro cúbico, é igual a
    %                         \begin{enumerate}
    %                             \item 3750
    %                             \item 1500
    %                             \item 1250
    %                             \item 375
    %                             \item 150
    %                         \end{enumerate}
    %                     \end{column}
    %                 \end{columns}
    %             \end{enumerate}
    %         \end{column}
    %     \end{columns}
    % \end{frame}
    
    % \begin{frame} %Questão 04-B
    %     \frametitle{Geometria}
    %     \begin{columns}[T, totalwidth=.9\textwidth]
    %         \begin{column}{\textwidth}
    %             \begin{enumerate}
    %                 \vest{ESPM-2018} O número de anagramas da palavra \textbf{COLEGA} em que as letras L, E e G aparecem juntas em qualquer ordem é igual a:
                    
    %                 \vspace{1.6em}

    %                 \centering
    %                 \begin{minipage}{.8\textwidth}
    %                     \begin{enumerate*}[itemjoin={\hfill}]
    %                             \item 72
    %                             \item 144
    %                             \item 120
    %                             \item 60
    %                             \item 24
    %                 \end{enumerate*}
    %                 \end{minipage}
    %             \end{enumerate}
    %         \end{column}
    %     \end{columns}
    % \end{frame}

    % \begin{frame} %Questão 05-A
    %     \frametitle{Geometria}
    %     \begin{columns}[T, totalwidth=.9\textwidth]
    %         \begin{column}{\textwidth}
    %             \begin{enumerate}
    %                 \vest{ENEM-2017} Suponha que para um trem trafegar de uma cidade à outra seja necessária a construção de um túnel com altura e largura iguais a $10$m. Por questões relacionadas ao tipo de solo a ser escavado, o túnel deverá ser tal que qualquer seção transversal seja o arco de uma determinada parábola, como apresentado na Figura 1.

    %                 Deseja-se saber qual a equação da parábola que contém esse arco. Considere um plano cartesiano com centro no ponto médio da base da abertura do túnel, conforme Figura 2.

    %                 A equação que descreve a parábola é:
    %                 \vspace{1em}
    %                 \begin{columns}[T, totalwidth=.9\textwidth]
    %                     \begin{column}{.3\textwidth}
    %                         \includegraphics[width=\textwidth]{monitoria_questao05_1.png}
    %                     \end{column}
    %                     \begin{column}{.3\textwidth}
    %                         \includegraphics[width=\textwidth]{monitoria_questao05_2.png}
    %                     \end{column}
    %                     \begin{column}{.2\textwidth}
    %                         \begin{enumerate}
    %                             \item $\displaystyle y = - \frac{2}{5}x^2 +10$\vspace{0.5em}
    %                             \item $\displaystyle y = \frac{2}{5}x^2 +10$\vspace{0.5em}
    %                             \item $y=-x^2 +10$\vspace{0.5em}
    %                             \item $y=x^2 -25$\vspace{0.5em}
    %                             \item $y=-x^2+25$
    %                         \end{enumerate}
    %                     \end{column}
    %                 \end{columns}
    %             \end{enumerate}
    %         \end{column}
    %     \end{columns}
    % \end{frame}

    % \begin{frame} %Questãao 06-B
    %     \frametitle{Geometria}
    %     \begin{columns}[T, totalwidth=.9\textwidth]
    %         \begin{column}{\textwidth}
    %             \begin{enumerate}
    %                 \vest{PUCR-RS/2013} Um desafio matemático construído pelos alunos do Curso de Matemática tem as peças no formato de um cone. A figura abaixo representa a planificação de uma das peças construídas.
    %                 \vspace{0.5em}
    %                 \begin{columns}[T, totalwidth=.9\textwidth]
    %                     \begin{column}{.4\textwidth}
    %                         \centering
    %                         \includegraphics[width=.9\textwidth]{lista_q06.png}
    %                     \end{column}
    %                     \begin{column}{.4\textwidth}
    %                         A área total dessa peça é de \underline{\hspace{.5cm}} $\text{cm}^2$

    %                         \vspace{1em}
                            
    %                         \begin{enumerate*}[itemjoin={\hfill}]
    %                             \item $10\pi$
    %                             \item $16\pi$
    %                             \item $20\pi$
    %                             \item $28\pi$
    %                             \item $40\pi$
    %                             % gabarito B 
    %                         \end{enumerate*}
    %                     \end{column}
    %                 \end{columns}
    %             \end{enumerate}
    %         \end{column}
    %     \end{columns}
    % \end{frame}

\end{document}